\documentclass[9pt,twocolumn]{article}
\usepackage[utf8]{inputenc}
\usepackage[margin=0.40in,top=0.38in,bottom=0.38in,columnsep=0.20in]{geometry}
\usepackage{amsmath,amssymb,amsfonts}
\usepackage{booktabs,tabularx,array,multirow}
\usepackage{microtype}
\usepackage{xcolor}
\usepackage{titlesec}
\usepackage{hyperref}

\hypersetup{
    colorlinks=true,
    linkcolor=blue!70!black,
    citecolor=blue!70!black,
    urlcolor=blue!70!black
}

% Professional compact section formatting
\titleformat{\section}{\bfseries\fontsize{9.5pt}{11pt}\selectfont\color{blue!80!black}}{\thesection}{0.4em}{}[\titlerule]
\titlespacing*{\section}{0pt}{4pt plus 1pt minus 1pt}{2.5pt plus 1pt minus 1pt}

\titleformat{\subsection}{\bfseries\fontsize{8.5pt}{10pt}\selectfont}{\thesubsection}{0.3em}{}
\titlespacing*{\subsection}{0pt}{3pt plus 1pt minus 1pt}{1.5pt plus 1pt minus 1pt}

\setlength{\parindent}{0pt}
\setlength{\parskip}{2.0pt plus 0.5pt minus 0.5pt}

\begin{document}

\twocolumn[
  \begin{center}
    {\Large\bfseries The Relay: Resilient Research Paper Assistant\par}
    \vspace{2pt}
    {\footnotesize\textbf{Author:} Udbhaw Anand \quad|\quad \textbf{Repository:} \url{https://github.com/Udbhaw08/the-relay} \quad|\quad \textbf{Deliverable Run ID:} \texttt{phase9\_final\_verified} (Seed 42, 125 Requests)\par}
  \end{center}
  \vspace{-5pt}
  \hrule
  \vspace{3pt}
]

\section{1. What I Built and Why}
I built \textbf{The Relay}, an agentic research assistant in Python answering technical engineering queries over three arXiv papers: \textit{Attention Is All You Need} (1706.03762v7), \textit{Moshi: A Speech-Text Foundation Model} (2410.00037v1), and \textit{Dynamic Chunking for H-Net} (2507.07955v2).

The system executes four decoupled stages:
\begin{enumerate}
    \setlength{\itemsep}{0pt}\setlength{\parskip}{0pt}
    \item \textbf{Query Parser \& Intent Guard:} Extracts technical keywords, identifies scope, and detects unanswerable questions early.
    \item \textbf{Hybrid Passage Retrieval:} Combines BM25Okapi lexical matching with dense vector search (\texttt{all-MiniLM-L6-v2}), fused via Reciprocal Rank Fusion ($\mathrm{RRF}(d) = \sum \frac{1}{60 + r(d)}$), followed by Cross-Encoder re-ranking.
    \item \textbf{Synthesis \& Citation Drafter:} Enforces strict bracketed anchors \texttt{[Paper, Page X, Section Y]}.
    \item \textbf{Citation Verification Guard:} Verifies claims against chunk provenance, triggering a single query-expansion loop-back on ungrounded claims and refusing unanswerable queries.
\end{enumerate}

\textbf{The One Decision That Shaped It Most:} 
Enforcing \textbf{zero silent fallback with adaptive retrieval degradation}. Upstream provider limits are never masked by silently routing queries to unannounced third-party models. Instead, when rate limits or queue pressures mount, the orchestrator adaptively throttles retrieval from Setting B/C down to compact Setting A (reducing prompt context tokens by $\sim$70\%), preserving service uptime while ensuring complete telemetry provenance.

\section{2. The Real Service and My Stand-In}
The system targets the \textbf{GroqCloud API} (\texttt{api.groq.com/openai/v1}) using model \texttt{openai/gpt-oss-20b} (following the deprecation of legacy \texttt{llama-3.3-70b-versatile} on August 16, 2026). All benchmark evaluations were executed using our calibrated \textbf{Stand-in Simulator} because the friend's account authorization and Free-tier status remain unverified.

\begin{table}[h]
\centering
\caption{\footnotesize Table 1: Published Groq Free-plan limits (documentation; not observed on key owner's account)}
\vspace{2pt}
\fontsize{7.2pt}{8.5pt}\selectfont
\begin{tabularx}{\columnwidth}{l X l c}
\toprule
\textbf{Limit} & \textbf{Published Value} & \textbf{Source / How Found} & \textbf{Conf.} \\
\midrule
RPM & 30 RPM (Free plan) & Official Groq Docs & High \\
TPM & 8,000 TPM & Groq Docs / \texttt{x-ratelimit-limit-tokens} & High \\
RPD & 1,000 RPD & Groq Docs / \texttt{x-ratelimit-limit-requests} & High \\
TPD & 200,000 TPD & Official Groq Docs & High \\
In-Flight & Dynamic (1--5 conc.) & Groq Cloud Documentation & Med \\
Hit Code & HTTP 429 & Groq API Error Specification & High \\
Recovery & Dynamic & Headers (\texttt{retry-after}, reset) & High \\
\bottomrule
\end{tabularx}
\end{table}

\textbf{Header Semantics \& Safety Caps:} Groq officially documents \texttt{x-ratelimit-limit-requests} as Requests Per Day (RPD), NOT RPM. \texttt{x-ratelimit-limit-tokens} corresponds to TPM. Local safety caps default conservatively (30 RPM, 8K TPM, 1K RPD, 200K TPD) to protect shared organization quotas. The friend's organization Limits page remains the source of truth for account allocations.

\textbf{Stand-in Methodology:} Latency is generated via a truncated Gaussian distribution ($\mu = 1200\,\mathrm{ms}$, $\sigma = 300\,\mathrm{ms}$, $\min = 150\,\mathrm{ms}$) with a token generation rate of $15\,\mathrm{ms/token}$ ($\approx 66\,\mathrm{tok/s}$). Contention is modeled via:
$$\mathrm{TTFT}_{\mathrm{eff}} = \mathrm{TTFT}_{\mathrm{base}} \times \left(1 + 0.08\,N_{\mathrm{in\text{-}flight}}^{1.15}\right)$$
TPM and TPD limits are enforced via real tracked rolling-window token consumption rather than manual booleans. As live Groq access was unverified, this is a synthetic model rather than an empirically calibrated trace; it cannot simulate Groq's hardware LPU SRAM compilation stalls or prompt prefix caching.

\section{3. Results}
All 125 benchmark requests were executed offline under run ID \texttt{phase9\_final\_verified} (seed 42), logging to \texttt{data/runs/phase9\_final\_verified/results.csv}.

\begin{table*}[t]
\centering
\caption{\footnotesize Table 2a: Latency Across Load Scenarios \quad|\quad Table 2b: Resilience Outcomes (Run ID: \texttt{phase9\_final\_verified})}
\vspace{2pt}
\fontsize{7.2pt}{8.5pt}\selectfont
\begin{tabular*}{\textwidth}{@{\extracolsep{\fill}} l c c c c c c @{}}
\toprule
\textbf{Scenario} & \textbf{Type} & \textbf{Reqs} & \textbf{Duration (s)} & \textbf{Offered Rate} & \shortstack{\textbf{Simulated TTFT}\\\textbf{p50 / p95 / p99 (ms)}} & \shortstack{\textbf{Total Latency}\\\textbf{p50 / p95 / p99 (ms)}} \\
\midrule
steady & stand-in & 10 & 70.9s (17.88s disp) & 30.2 req/min (int 1.99s) & 1,840.4 / 2,132.8 / 2,137.5 & 13,519.8 / 16,109.8 / 16,239.5 \\
burst & stand-in & 20 & 66.2s (59.57s disp) & Burst (20 concurrent) & 2,630.0 / 3,760.0 / 3,983.5 & 6,030.1 / 58,282.0 / 58,680.8 \\
sustained\_overload & stand-in & 25 & 120.1s (115.64s disp) & 13.0 req/min & 1,515.8 / 2,334.6 / 2,389.2 & 5,300.9 / 54,717.5 / 54,978.1 \\
per\_minute\_halved & stand-in & 15 & 236.7s (194.12s disp) & 4.6 req/min & 1,393.0 / 1,666.8 / 1,700.4 & 3,958.4 / 23,743.8 / 38,788.2 \\
daily\_exhausted & stand-in & 10 & 2.0s & 306.1 req/min & N/A & N/A \\
\bottomrule
\end{tabular*}

\vspace{3pt}

\begin{tabular*}{\textwidth}{@{\extracolsep{\fill}} l c c c c c c c @{}}
\toprule
\textbf{Scenario} & \textbf{Succeeded \%} & \textbf{Rate-Limited \%} & \textbf{Timed Out \%} & \textbf{Retries/Req} & \textbf{Peak Queue} & \textbf{Tokens/Ans} & \textbf{Cost / 1,000 Answers} \\
\midrule
steady & 70.0\% & 30.0\% & 0.0\% & 0.30 & 5 & 41.1 & N/A — no live provider usage \\
burst & 65.0\% & 25.0\% & 10.0\% & 0.95 & 20 & 36.1 & N/A — no live provider usage \\
sustained\_overload & 40.0\% & 36.0\% & 24.0\% & 0.96 & 15 & 33.1 & N/A — no live provider usage \\
per\_minute\_halved & 73.3\% & 20.0\% & 6.7\% & 0.60 & 0 & 37.0 & N/A — no live provider usage \\
daily\_exhausted & 0.0\% & 100.0\% & 0.0\% & 0.00 & 0 & 0.0 & N/A — no live provider usage \\
\bottomrule
\end{tabular*}
\vspace{-3pt}
\end{table*}

\begin{table*}[t]
\centering
\caption{\footnotesize Table 3: Retrieval Settings Evaluation and Factual Accuracy Breakdown}
\vspace{2pt}
\fontsize{7.2pt}{8.5pt}\selectfont
\begin{tabular*}{\textwidth}{@{\extracolsep{\fill}} l c c c c c c c @{}}
\toprule
\shortstack{\textbf{Setting}\\\textbf{(Chunk, Fetch, Rerank)}} & \textbf{Questions} & \shortstack{\textbf{Correct \&}\\\textbf{Cited \%}} & \textbf{Wrong \%} & \shortstack{\textbf{Correctly}\\\textbf{Refused \%}} & \textbf{Tokens/Ans} & \shortstack{\textbf{Simulated TTFT}\\\textbf{p50/p95 (ms)}} & \shortstack{\textbf{Total Time}\\\textbf{p50/p95 (ms)}} \\
\midrule
Setting A: 250 tok, 3 fetch, no rerank & 15 & 8.3\% (1/12) & 91.7\% (11/12) & 100.0\% (3/3) & 39.1 & 1,418.0 / 1,558.0 & 3,687.8 / 4,835.5 \\
Setting B: 500 tok, 5 fetch, rerank: yes & 15 & 8.3\% (1/12) & 91.7\% (11/12) & 100.0\% (3/3) & 37.2 & 1,418.0 / 1,566.8 & 3,773.9 / 5,125.3 \\
Setting C: 1000 tok, 7 fetch, rerank: yes & 15 & 8.3\% (1/12) & 91.7\% (11/12) & 100.0\% (3/3) & 33.1 & 1,418.0 / 1,573.3 & 3,617.9 / 4,767.6 \\
\bottomrule
\end{tabular*}

\vspace{3pt}

\begin{tabular*}{\textwidth}{@{\extracolsep{\fill}} l c c c c c c c c c @{}}
\toprule
\textbf{Setting} & \textbf{Ans Eval} & \textbf{Correct \& Cited} & \textbf{Partial} & \textbf{Wrong} & \textbf{Wrong Refused} & \textbf{Unans Eval} & \textbf{Correct Refused} & \textbf{Hallucinated} & \textbf{Infra Fails} \\
\midrule
setting\_a & 12/12 & 1 & 0 & 11 & 0 & 3/3 & 3 & 0 & 0/15 \\
setting\_b & 12/12 & 1 & 0 & 10 & 1 & 3/3 & 3 & 0 & 0/15 \\
setting\_c & 12/12 & 1 & 0 & 8 & 3 & 3/3 & 3 & 0 & 0/15 \\
\bottomrule
\end{tabular*}
\vspace{-3pt}
\end{table*}

\textbf{Metric Provenance:} Token metrics report estimated counts based on whitespace heuristics. Actual provider token fields are null because no live API calls were made. Cost is labeled \textit{"N/A — no live provider usage"}. First token latency is labeled as Simulated TTFT and is measured exclusively on streaming attempts (non-streaming requests set \texttt{first\_token\_ms = null}).

\textbf{Denominator \& Isolation Notice:} Pure retrieval benchmarks execute with adaptive degradation disabled (\texttt{allow\_degradation=False}) so Setting A cannot contaminate Settings B or C (15/15 pure requests per setting). Infrastructure failures are excluded from factual accuracy denominators and reported separately. Answerable accuracy (12 questions) and unanswerable refusal accuracy (3 questions) are segregated with exact sample counts.

\textbf{Chosen Default Setting:} \textbf{Setting B (500 tokens, 5 fetched, hybrid RRF re-ranking)} is the default. It balances semantic context depth against prompt token cost. Under rate-limit warnings, the engine dynamically degrades to \textbf{Setting A} (250 tokens, 3 fetched, BM25-only), shedding $\sim$70\% of context tokens.

\textbf{An Unexpected Number:} In steady load, bounded concurrent scheduling achieved exact 1.0s target start-to-start pacing (\textbf{mean measured interval: 1.0s}, \textbf{mean dispatch lateness: 6.66ms}) despite cold-start retrieval latencies reaching 13.5s. In Scenario 4, halving baseline RPM to 7 actively exercised rate limiting (20.0\% rate-limited, 73.3\% degraded success), confirming throttling is operational.

\section{4. Where the Time Goes}
Profiling slow request \texttt{steady\_stand-in\_001} (Total: \textbf{15,731.4 ms}):
\begin{itemize}
    \setlength{\itemsep}{0pt}\setlength{\parskip}{0pt}
    \item \textbf{Queue Wait Time:} 0.3 ms (admitted immediately).
    \item \textbf{Retrieval Time:} 12,831.8 ms (\textbf{81.6\% of duration}). Dense embedding cold-start across all corpus chunk partitions.
    \item \textbf{Inference Time:} 2,899.2 ms (18.4\% of duration: simulated TTFT 2,138.7 ms, generation 760.5 ms across 46 tokens).
    \item \textbf{Verification Time:} 0.3 ms ($<0.01\%$: regex extraction and page/anchor verification).
    \item \textbf{Warm Comparison:} Once vector weights are memory-resident (\texttt{steady\_stand-in\_009}), retrieval drops to \textbf{10.6 ms} (99.9\% reduction), reducing total request latency to 2,047.9 ms.
\end{itemize}

\section{5. When a Limit Is Hit}
\textbf{Per-Minute Limit (RPM/TPM):} Handled via an in-memory queue. Requests execute truncated exponential backoff with full jitter ($T \sim \mathcal{U}(0, \min(M, B \times 2^k))$) and adaptively throttle retrieval to Setting A (250 tokens, BM25-only), shedding $\sim$70\% of prompt token load.

\textbf{Per-Day Limit (RPD/TPD):} Stops retries immediately (0 retries). Requests fail fast ($<1\mathrm{ms}$) with an explicit notice stating dynamic reset estimation, preventing thundering-herd queue lockup.

\section{6. How This Library Will Mislead a Naive Retriever}
\begin{enumerate}
    \setlength{\itemsep}{0pt}\setlength{\parskip}{0pt}
    \item \textbf{Lexical Collision on ``Attention'':} Queries regarding \textit{Moshi's} temporal stream masking match Vaswani's translation masking in \textit{Attention Is All You Need}, returning translation equations instead of audio synchronization.
    \item \textbf{Term Collision on ``Chunking'':} Queries on dynamic chunking benefits over tokenizers match \textit{Moshi's} 80ms acoustic frame stride rather than \textit{H-Net's} byte-level hierarchy.
    \item \textbf{Vocabulary Mismatch on ``Tokens'':} Asking about discrete audio speech tokens retrieves BPE subwords from Vaswani et al. rather than Mimi neural codec RVQ codebooks.
\end{enumerate}

\section{7. Retries}
\textbf{Policy:} Truncated exponential backoff with true full jitter ($T \sim \mathcal{U}(0, \min(30.0, 2.0 \times 2^k))$), maximum 2 retries, respecting provider \texttt{Retry-After} headers up to a 60s safety cap.

\textbf{Where Retrying Made Things Worse:} Retrying on \textbf{daily budget exhaustion} caused catastrophic queue saturation. Classifying daily exhaustion as a non-retryable error eliminated retry loops completely.

\section{8. What the Person Sees}
When daily quota is depleted, the user receives an explicit, transparent notice:
\begin{quote}
\footnotesize
\textit{``[Notice] Daily Query Budget Exhausted: The system's daily LLM inference capacity has reached its limit for today. The next quota allocation resets dynamically (or when upstream headers indicate). Your request was rejected locally to protect shared quotas. Please consult the local paper index or retry once the quota window refreshes.''}
\end{quote}

\vspace{-2pt}
\hrule
\vspace{2pt}
{\tiny\textit{Note on Extended Analysis:} Comprehensive phase execution matrices, architectural diagrams, extended reflections, feedback on Vaani, and assistant collaboration logs are documented in \url{docs/implementation/FINAL_REPORT.md}.}

\end{document}
